% !TeX root = ../../Thesis.tex
\ThesisChapter{\ThesisText{Introduction}{Einleitung}}{ch:introduction}
  {TODO: Add a brief chapter overview. \\
  This chapter motivates the topic, states the research questions and objectives, describes the methodology and contributions, and outlines the structure of the thesis.}

This chapter is deliberately written as continuous prose, with no section headings.
Both the research questions and the contributions can be set as lists, and that is what makes them findable, so headings are not needed.
As a general guideline, the introduction is three to four pages long.

% ---------------------------------------------------------------------------
% (1) Motivation and problem statement
% ---------------------------------------------------------------------------
% TODO: Introduce the topic and explain the motivation and problem addressed by the thesis.

Provide enough context for readers who are not yet familiar with the specific subject.
Explain why the topic is relevant, which practical or scientific challenges exist, and why the problem is worth investigating.

Clearly define the problem addressed by the thesis and briefly indicate the limitations of existing solutions.
This part should lead to the research questions without already describing the proposed solution in detail.

As a general guideline, this part should be approximately one page long.


% ---------------------------------------------------------------------------
% (2) Research questions
% ---------------------------------------------------------------------------
% TODO: Formulate the research questions of the thesis.

The research questions follow directly from the problem stated above.
They define what this thesis investigates, and every one of them must be answered by the end of it.

As a general guideline, three to six clearly formulated questions are recommended.
They do not have to be phrased as questions; concise declarative statements work equally well.

Each item below is numbered RQ1, RQ2, and so on, and carries a label.
Referencing that label anywhere else in the thesis prints the same name, which is how the conclusion answers each question by name.

\begin{ThesisResearchQuestions}
  \item\label{rq:example-one} TODO: Replace with the first research question.
  \item\label{rq:example-two} TODO: Replace with the second research question.
  \item\label{rq:example-three} TODO: Replace with the third research question.
\end{ThesisResearchQuestions}

% ---------------------------------------------------------------------------
% (3) Objectives
% ---------------------------------------------------------------------------
% TODO: Derive the objectives of the thesis from the research questions above.
%
A research question states what is asked.
An objective states what will be done in order to answer it.

Make the link explicit, so that no objective stands on its own and no research question is left without one.
Keep the objectives concrete enough that it can be judged at the end whether they were achieved.

% ---------------------------------------------------------------------------
% (4) Methodology
% ---------------------------------------------------------------------------
% TODO: Describe how the research questions will be answered.

If a specific research framework or methodology is used, introduce it here and relate it to the stages of the thesis.
Otherwise describe the approach in practical terms: what will be designed, built, measured, and compared.

State the criteria by which the outcome will be judged, so that the evaluation chapter has something to measure against.

% ---------------------------------------------------------------------------
% (5) Contributions
% ---------------------------------------------------------------------------
% TODO: State what this thesis actually produced.

A contribution is a result, not an activity.
Implementing a scheduler is a task; showing that work stealing reduces tail latency by a third on up to 64 nodes is a contribution.

As a general guideline, state three to five contributions.
Name for each one the chapter that delivers it, so that a reader can go straight to the evidence.
Typical contributions include a concept or design, an implementation that makes it usable, an evaluation that produces new evidence, and a publicly available artifact.

 Keep this list consistent with the abstract, which states the same contributions in prose.

\begin{ThesisContribution}
  \item\label{c:example-one} TODO: Replace with the first contribution.
  \item\label{c:example-two} TODO: Replace with the second contribution.
  \item\label{c:example-three} TODO: Replace with the third contribution.
\end{ThesisContribution}


% ---------------------------------------------------------------------------
% (6) Structure of the thesis
% ---------------------------------------------------------------------------
% TODO: Give a brief overview of the remaining chapters.

Explain how each chapter contributes to the objectives and research questions.
This overview may be written as continuous text, as a structured list, or with the help of a figure.

It is often useful to indicate which chapter addresses which research question.
